\documentclass[11pt]{article}

\usepackage[margin=0.9in]{geometry}
\usepackage[T1]{fontenc}
\usepackage[utf8]{inputenc}
\usepackage{amsmath,amssymb}
\usepackage{booktabs,longtable,tabularx,array}
\usepackage{enumitem}
\usepackage{xcolor}
\usepackage[hidelinks]{hyperref}

\newcommand{\code}[1]{\texttt{\detokenize{#1}}}
\newcommand{\rhoctl}{\rho}
\newcommand{\kappactl}{\kappa}
\newcommand{\Kbase}{K_{\mathrm{base}}}
\newcommand{\Keff}{K_{\mathrm{eff}}}
\newcommand{\Kref}{K_{\mathrm{ref}}}
\newcommand{\Kmax}{K_{\max}}
\newcommand{\reff}{r_{\mathrm{eff}}}
\newcommand{\ssafe}{s_{\mathrm{safe}}}
\newcommand{\truefamily}{\mathrm{true\_family}}

\setlength{\parskip}{0.35em}
\setlength{\parindent}{0pt}
\setlist[itemize]{leftmargin=1.5em,itemsep=2pt,topsep=3pt}
\setlist[enumerate]{leftmargin=1.7em,itemsep=2pt,topsep=3pt}
\sloppy

\title{\textbf{Problem and Fix Brief}\\
\large Hidden ecological parameters $(r,K)$ are leaked to the agents}
\author{Readable TeX version of \code{SERVER_FIX_BRIEF_hide_rK.md}}
\date{2026-07-15}

\begin{document}
\maketitle
\tableofcontents
\newpage

\section{Purpose}

This brief is for the code-aware server agent. The current real-ecology
experiment exposes ecological parameters that the intended uncertainty setting
should hide. The server agent should verify the exposure against the current
code, then implement the hidden-\(r,K\) fix while preserving the simulator's
private use of the real ecological tables.

This TeX file is a readable companion to \code{SERVER_FIX_BRIEF_hide_rK.md}; the
Markdown brief remains the implementation source of record.

\section{First Deliverable: Root-Cause Note}

Before implementing anything, the server agent should trace where exposure of
\(r\), \(K\), and true-family metadata was decided. The goal is to identify the
source of the error, not merely patch symptoms.

\begin{enumerate}
  \item \textbf{Which document decided it.} Check the design/spec docs,
  especially implementation-plan E1 and E8:
  \begin{itemize}
    \item E1 exposes population identity and \((\reff,\Keff)\) as observed
    inputs.
    \item E1 also treats the \(\lambda\) profile, \(\Kbase\), and caps as
    observed inputs.
    \item E8 says the methods only need to read already-known
    \((\reff,\Keff)\).
    \item The algorithm/feature spec includes public-control terms in the
    feature vector.
  \end{itemize}
  Quote the exact lines in the root-cause note.

  \item \textbf{Which code enforces it.} Check the environment's observed-input
  construction, the public offline tuple, the feature builders for
  \code{refplan}/\code{bamcts}/\code{ogsrl}, and the native-solver table reads.
  Give file, function, line references, and git blame.

  \item \textbf{Spec error or code deviation.} State whether this was a flawed
  instruction faithfully implemented or a deviation from the docs. Working
  hypothesis: this is a design/spec error propagated to code, not a stray coding
  bug.

  \item \textbf{Why it was missed.} There was a reward-leakage guard, but no
  parameter-leakage guard and no explicit public/private schema derived from the
  motivation.

  \item \textbf{Prevention rule.} Every observed input should pass this test:
  does exposing it collapse the uncertainty the experiment is meant to test?
  Also: identity is not demographics.
\end{enumerate}

\section{The Problem}

The ecological parameters that were meant to be hidden--the per-action growth
rate \(r(a)\) / \(\reff\) and the per-population carrying capacity
\(\Kbase/\Keff\)--are exposed to the agents as observed inputs. They are
published in the action/species tables and fed to method-facing data, features,
belief caches, planners, or native solver construction.

This turns the current experiment into a structure-known or public-demographics
setting, not the intended hidden-parameter uncertainty setting. The completed
motivation-native run is still a valid result: it shows that native mechanistic
solvers are strong when demographics/action effects are public. But it should be
reported as a public-demographics / parameter-known data-rich regime, not as the
hidden-\(r,K\) experiment. Here, ``data-rich'' means the demographics/action
effects are well-characterized and public, not that the transition buffer is
large.

\section{Where the Leak Is}

The server usage audit has already confirmed the broad exposure pattern; the
server agent should re-verify these paths at fix time so implementation follows
the current code.

\begin{itemize}
  \item \textbf{Spec exposure.} Implementation-plan E1/E8 exposes
  \((\reff,\Keff)\), \(\Kbase\), caps, and the \(\lambda\) profile as known
  observed context.

  \item \textbf{Public offline tuple.} The public dataset includes
  \(\rho_{t+1}\), which is the set-point \(\reff\) when \(d_r=1\), and
  \(\Keff_{t+1}\).

  \item \textbf{General learner features.} Public-control terms such as
  \(\rho_{t+1}\), \(\kappa_{t+1}/\Kref\), and \(\Keff/\Kref\) enter learned
  dynamics or policy features.

  \item \textbf{Native solvers.} \code{moor_native} reads action-specific growth
  set-points and species-capacity structure to build a table-driven mechanistic
  model. It is near-exact on Ricker and misspecified but still strong off-Ricker.
  \code{plus_native} reads the same public structure and updates mainly a
  posterior over functional form.
\end{itemize}

Useful grep handles for the server code:
\[
\code{r_eff},\quad \code{delta_r},\quad \code{r_setpoint},\quad
\code{K_base},\quad \code{K_eff},\quad \code{action_effects},\quad
\code{species.csv}.
\]
Also inspect the dynamics-model feature construction for
\code{refplan}/\code{bamcts}/\code{ogsrl} and model construction for
\code{moor_native}/\code{plus_native}.

\section{Not a Duplicate of Prior Runs}

The earlier \code{psafe_overnight_*} run and the later
\code{motivation_native_*} run both used the same public-table environment. They
differed in ecological-solver strength: adapted particle-MPC baselines versus
native tabular solvers. Neither run hid \(r\) or \(K\).

Therefore the hidden-\(r,K\) redesign is new work. It is not a duplicate of
last week's setting.

\section{Why It Matters}

\begin{itemize}
  \item It removes the parameter-uncertainty layer the paper claims to test.
  \item It gives native solvers table-derived mechanistic models, so the
  ``native beats general'' result is valid for a public-demographics /
  data-rich comparison but not evidence about the intended hidden-parameter
  regime.
  \item The general learners appear model-limited partly because they compete
  against native solvers that were handed exact demographic/action structure.
  \item The current run should be kept as a labelled data-rich comparison arm,
  not discarded.
\end{itemize}

\section{Current vs Corrected Method-Facing Information Flow}

The left column describes what the current public-demographics / parameter-known
run allows. The right column describes what the corrected hidden-demographics /
structure-unknown run must enforce. The simulator/evaluator may still use real
tables privately; methods may not.

Here, structure-unknown means both hidden \(r/K\) and action-effect parameters on
the parameter axis, plus private \(\truefamily\) on the model-form axis. The
schema fix applies to every method in the codebase, including \code{mopo},
\code{delphic}, and adapted \code{moor}/\code{plus}. The headline hidden-\(r,K\)
run uses five methods: \code{refplan}, \code{bamcts}, \code{ogsrl},
\code{moor_native}, and \code{plus_native}.

\begin{longtable}{p{0.18\linewidth}p{0.38\linewidth}p{0.38\linewidth}}
\toprule
Scope / method & Current run: public-demographics / parameter-known & Corrected hidden-\(r,K\) run \\
\midrule
\endfirsthead
\toprule
Scope / method & Current run: public-demographics / parameter-known & Corrected hidden-\(r,K\) run \\
\midrule
\endhead
General public schema &
Public dataset, belief cache, and configs may expose \(\rho\), \(\kappa\),
\(\Keff\), \(\Kbase\), \(\Kref\), \(\Kmax\), action set-points, \(\Delta K\),
\(\Delta N\), K-derived \(\ssafe\), population table lookups, and family/config
identity. &
Methods may consume only sanitized public data: categorical \code{pop_id} if
chosen, \(o_t\), \(a_t\), \(\mathrm{cost}(a_t)\), \(R_t\) as label only,
\(o_{t+1}\), and \code{done}. All \(r/K\) aliases, table effects, K-derived
safety/normalization terms, and true-family metadata stay private. \\

\code{refplan} &
Learns dynamics with \(\rho\), \(\kappa/\Kref\), and \(\Keff/\Kref\); MPC
carries public controls and uses \(\Kref\)/safety reward terms. &
Fit dynamics and plan from sanitized belief/observation/action features only.
If action effects matter, infer them from data. Do not pass exact \(\rho\),
\(\kappa\), \(\Keff\), \(\Kref\), or K-derived safety features. \\

\code{bamcts} &
Uses the same control-aware learned dynamics as \code{refplan}; tree keys and
rollouts condition on \(\rho\) and \(\kappa\). &
Tree state and rollouts must be indistinguishable for histories with the same
\((o,a)\) but different private \(r/K\). Learned dynamics cannot receive
public-control aliases or exact table-derived action effects. \\

\code{ogsrl} &
Uses \(\rho\), \(\kappa/\Kref\), \(\Keff/\Kref\), \(\ssafe\), and unsafe/safety
features in dynamics, actor features, guardian features, rollouts, and
feasibility checks. &
Remove \(r/K\)-derived features from dynamics, actor, guardian, rollout, and
safety inputs. Safety-aware behavior must learn risk from public trajectories or
use a public threshold not derived from hidden \(\Kbase\). \\

\code{mopo} &
Uses the same learned-dynamics public-control channel as the other general
planners; planning reward uses \(\Kref\) and safety configuration. &
Use sanitized features only. Model uncertainty should include uncertainty about
hidden action effects and capacity rather than conditioning on exact public
controls. \\

\code{delphic} &
Does not build an ecological transition model, but receives shared belief
features containing \(\rho\), \(\kappa/\Kref\), \(\Keff/\Kref\), and safety
terms. &
Build compatible-world and Q features from sanitized public beliefs only.
Remove hidden-parameter aliases from the shared feature vector so
\code{delphic} cannot inherit the leak indirectly. \\

adapted \code{moor} &
Uses exact action set-points, \(\Delta K\), \(\Delta N\), and species-capacity
structure in a mechanistic Ricker proposal; mainly fits a K-like scale in
set-point mode. &
Mechanistic proposal cannot call table-derived exact action effects or species
capacity. It must fit required Ricker parameters/effects from public data, or be
reported only in the public-demographics comparison arm. \\

adapted \code{plus} &
Uses exact table-derived action effects and species-capacity ranges inside a
candidate mechanistic bank; candidate axis is mostly \(K\) because growth
set-points are public. &
Candidate parameters and priors must be derived from sanitized public data, not
from species/action tables. If retained, report separately from the naive
headline if it performs closed-world model selection. \\

\code{moor_native} &
Fits \(\hat K\), but reads exact action growth set-points, \(\Delta K\),
\(\Delta N\), costs, \(\Kref\), and \(\ssafe\) to build its native Ricker
solver. &
Headline naive baseline. Fit/estimate Ricker parameters and action effects from
sanitized public data, then solve a single-Ricker model. Do not read
species/action tables or K-derived safety/normalization values. \\

\code{plus_native} &
Builds one native solver per candidate family using exact table-derived \(r/K\),
action effects, reward, and safety structure; updates mainly the family
posterior. &
Separate closed-world model-selection baseline. It may keep a family posterior,
but all parameter candidates/priors must be fit from sanitized public data. Do
not redesign open-world PLUS in this first hidden-\(r,K\) repair. \\
\bottomrule
\end{longtable}

\section{Recommended First-Fix Decisions}

Unless the PI overrides them:

\begin{enumerate}
  \item Use known population label, unknown demographics. \code{pop_id} is a
  categorical token only, never a table lookup path.
  \item Use \code{expose_rk=hidden} as the intended default, and keep
  \code{expose_rk=full} to reproduce the current public-demographics /
  data-rich comparison arm. Keep both regimes runnable behind one interface; do
  not delete or overwrite \code{full}.
  \item Do not add a noisy-\(r/K\) setting until the hidden setting is stable.
  \item Keep \(\Kref\) private/evaluator-only.
  \item Keep any \(\ssafe\) derived from \(\Kbase\) private/evaluator-only.
  \item If a public safety threshold is needed later, define it independently of
  hidden \(\Kbase\).
  \item Use \code{moor_native} as the headline naive ecological baseline.
  \item Report \code{plus_native} separately as closed-world form selection.
  \item Postpone open-world PLUS.
\end{enumerate}

\section{The Fix: Hide \(r\) and \(K\) from All Agents}

Goal: move from the structure-known regime to the structure-unknown regime,
where per-action growth and per-population carrying capacity are latent and must
be inferred from offline data by every method.

\subsection{Stop Feeding \(\reff\) / \(K\) to Agent Inputs}

\begin{itemize}
  \item Remove \(\rho_{t+1}\), \(\Keff\), and \(\Kbase\)-derived quantities from
  general learner feature vectors and observed inputs passed to policy, belief,
  model, or planner code.
  \item The agent may know which action it took, \(a_t\), but not the growth
  rate that action induces and not the population's \(K\).
  \item Keep action cost observed, because cost does not reveal \(r\) or \(K\)
  by itself. Hide the action-to-\((r,K)\)-effect mapping.
\end{itemize}

\subsection{Make Native Solvers Estimate \(r\) and \(K\)}

This is the crux. If only general methods lose \(r/K\) while
\code{moor_native}/\code{plus_native} keep table access, the asymmetry gets
worse.

\begin{itemize}
  \item \code{moor_native} must fit its Ricker \(r\), \(K\), and action effects
  from sanitized public data, not from tables.
  \item \code{plus_native} must estimate \(r\) and \(K\) from data while keeping
  its existing form-candidate set for this first fix.
  \item For \code{plus_native}, per-candidate parameters/priors must be fit from
  sanitized data. Do not add explicit \(r/K\) candidate axes, distractor forms,
  or open-world candidates in this first repair.
  \item Add an explicit \code{fit_from_public_data()} entry point for native
  solvers. It may take only sanitized public data and belief/observation
  features derived from that data. It may not take or derive values from
  \code{species.csv}, \code{action_effects_long.csv}, \code{species_lambda.csv},
  or any true \(r/K\) table.
\end{itemize}

\subsection{Keep Tables Private}

The simulator/evaluator may continue to use real \(r/K\) tables and sidecars.
Methods must not receive them directly or through aliases, configs, filenames,
features, belief caches, reward decomposition, or table lookup.

\subsection{Population Identity}

Recommended setting: known species label, unknown demographics. The agent may
receive \code{pop_id} as a categorical label, but no method may use
\code{pop_id -> species.csv} to recover \(r\), \(K\), caps, \(\Kref\), or
\(\ssafe\).

Fully hidden population identity is a harder later variant, not the default
first repair.

\subsection{Explicit Regime Split: Keep Both Full and Hidden}

Implement the repair as an explicit regime split, not as a global overwrite of
the current \(r,K\)-known behaviour. Route both regimes through one shared
interface, for example \code{expose_rk \(\in\) \{full, hidden\}} or
\code{info_regime \(\in\) \{public_demographics, hidden_demographics\}}.

\begin{itemize}
  \item \code{full}: keep the current public-demographics / parameter-known run
  as a labelled ablation or upper-bound arm. Modify this path only enough to
  label it and route it through the shared interface.
  \item \code{hidden}: the corrected hidden-demographics / structure-unknown
  setting. This is the intended default and headline uncertainty test.
  \item Record the regime in output directories, configs, manifests, and result
  summaries so the two arms cannot be confused.
  \item A noisy-\(r/K\) prior variant is optional later work only; it is not part
  of the required first repair.
\end{itemize}

\section{Public / Private Schema}

A side channel anywhere in this schema reopens the problem. Audit every method
input path against it.

\subsection{Hidden: Private Evaluator Sidecar Only}

Never expose the following to any method:

\begin{itemize}
  \item true abundance \(s_t\);
  \item structural latents \(C\), \(\theta\), and regime \(z\);
  \item all growth quantities and aliases: \(r(a)\), \(\reff\), set-point
  Ricker/LGM values, \(\rho_t\), \(\rho_{t+1}\), \(r_{\min}\), \(r_{\max}\),
  \(r_{\mathrm{base}}\), and per-population \(\lambda\) profiles;
  \item all capacity quantities and aliases: \(\Kbase\), \(\Keff_t\),
  \(\Keff_{t+1}\), \(\Kmax\), \(\Kref\), \(\kappa_t\), \(\kappa_{t+1}\),
  and \(\Delta K\);
  \item direct-state action magnitude: \(\Delta N\) / translocation magnitude;
  \item any table encoding those quantities: \code{species.csv},
  \code{action_effects_long.csv}, \code{species_lambda.csv}, and cap columns;
  \item the shared belief cache / feature builder if it contains or derives any
  hidden quantity. Sanitize at the source so no method, especially
  \code{delphic}, inherits the leak indirectly.
\end{itemize}

\subsection{Observed: Sanitized Public Dataset}

Methods may load only:
\[
(\code{pop_id},\; o_t,\; a_t,\; \mathrm{cost}(a_t),\; R_t,\; o_{t+1},\; \code{done})
\]
plus belief or observation features derived only from those fields. Here,
\code{pop_id} is a categorical label only, if the known-label design is used.
\(R_t\) is an offline learning label only.

\subsection{Indirect Capacity Leaks}

\(\Kref=\Kbase\) and \(\ssafe=c_{\mathrm{safe}}\Kbase\) leak the capacity scale.
Therefore:

\begin{itemize}
  \item \(\Kref\) should be private/evaluator-only.
  \item Any \(\ssafe\) derived from \(\Kbase\) should be private/evaluator-only.
  \item If a public safety threshold is needed later, define it independently of
  hidden \(\Kbase\).
\end{itemize}

\subsection{Reward Guard}

\(R_t\) may remain an offline learning label. None of the following may enter
policy, belief, model, or planner features:

\begin{itemize}
  \item reward decomposition;
  \item benefit term \(\alpha s'/(s'+\Kref)\);
  \item unsafe/collapse flag;
  \item true state \(s\);
  \item any \(s\)- or \(K\)-based diagnostic.
\end{itemize}

\section{Structural Uncertainty: Do Not Leak True Family}

This fix hides \(r\) and \(K\), but structural-form uncertainty must also be
preserved.

\begin{itemize}
  \item The true family and aliases--\code{true_family}, \code{family},
  \code{dynamics_family}, \code{model_family}, \code{map_family},
  \code{env.family}, and \code{cfg.kind}--are private evaluator/simulator
  metadata only.
  \item No method may receive true family as input, feature, config-conditioned
  model choice, table lookup, manifest shortcut, or filename/path-derived
  switch.
  \item \code{moor_native} is the headline naive baseline and remains a
  single-Ricker solver even when truth is Allee, theta-logistic, or
  regime-switching.
  \item \code{plus_native} is not the naive headline. If it keeps the candidate
  set \(\{\mathrm{Ricker}, \mathrm{Allee}, \theta\text{-logistic},
  \mathrm{regime}\}\), report it separately as closed-world model selection.
  \item Do not redesign \code{plus_native}'s candidate set in this first repair.
  Open-world PLUS, distractor forms, or truth-absent candidate sets are later
  experiments.
\end{itemize}

\section{Acceptance Tests}

\begin{enumerate}
  \item \textbf{No-r/K-access test.} Prove no method receives \(\reff\),
  \(\Kbase\), \(\Keff\), or any hidden alias as input or feature. Two rollouts
  with identical \((o,a)\) histories but different private \(r/K\) must be
  indistinguishable to the policy input.

  \item \textbf{Table-independence test.} Corrupt or shuffle
  \code{species.csv}, \code{action_effects_long.csv}, and
  \code{species_lambda.csv}. Every method's training, fit, belief features, and
  planning artifacts should be byte-identical if the public dataset is fixed.

  \item \textbf{Family-relabel test.} With the same public dataset but relabelled
  private \(\truefamily\), method training/fitting/planning artifacts must be
  unchanged. Only evaluator grouping/reporting may change.

  \item \textbf{Native solvers now estimate.} Confirm
  \code{moor_native}/\code{plus_native} fit \(r\) and \(K\) from data, not table
  reads. Their estimates should degrade under scarce/noisy data.

  \item \textbf{Identifiability sanity.} Check whether \(r\) and \(K\) are
  recoverable from the 4000-transition noisy-observation budget. If not, report
  it as the actual uncertainty, not a code bug.

  \item \textbf{Expected outcome, not a gate.} Native performance should drop
  relative to the leaked run because the native solvers are no longer near
  oracle. If it does not drop, report that as a result: \(r/K\) were identifiable
  from public data.

  \item \textbf{Regime split intact: both runs live.} \code{full} mode should
  reproduce the current method-facing fields and behaviour where reproducible;
  \code{hidden} mode must remove every hidden/private field listed above from
  all methods. The simulator/evaluator should use the private true tables in
  both regimes, and table-corruption or family-relabel tests should affect
  hidden-mode agents only through the public data, never through direct table
  lookup. Confirm the regime is recorded in every output/config/manifest.

  \item \textbf{Re-run headline comparison.} Run the hidden setting across all
  cells and report per-\(\sigma\), per-family, recoverable/sink split.

  \item Verify that all controlled-comparison settings remain unchanged: nine
  populations, four families, observation-noise sweep, reward modes, 4000
  transitions per cell, evaluation protocol, planner budget, true-state reward
  in evaluator, and private/public boundary.
\end{enumerate}

\section{Do Not Confuse This with Planning-Model-Source Ablation}

This fix changes the environment/observation information boundary: it stops
publishing \(r/K\) so demographics become latent for everyone.

The planning-model-source ablation is a separate experiment. It changes what
dynamics the general planners roll through on the current leaked environment.
That ablation is not a substitute for the hidden-\(r,K\) repair.

\section{Open PI Decisions}

\begin{enumerate}
  \item Population identity: known species label with unknown demographics
  (recommended) versus fully hidden population.
  \item Whether and when to add the noisy-\(r/K\) prior variant. Keep it off by
  default until \code{hidden} is stable. Keeping both \code{full} and
  \code{hidden} regimes behind one interface is already decided; it is a
  requirement, not an open decision.
  \item Whether and when to build open-world PLUS. Recommendation: later, not
  part of this first repair.
\end{enumerate}

\section{Deliverables}

The server agent should report back with:

\begin{enumerate}
  \item root-cause note: doc line, code path, git blame, whether spec error or
  code deviation, why it was missed, and one-line prevention rule;
  \item answers to all verify items and PI decisions with exact file/function
  paths;
  \item sanitized public dataset loader and private evaluator sidecar as a diff;
  \item \code{moor_native} and \code{plus_native} \code{fit_from_public_data()}
  implementations;
  \item all acceptance tests passing, especially table-independence,
  family-relabel, and the regime-split-intact test;
  \item corrected hidden-\(r,K\) run: best general versus \code{moor_native},
  Ricker versus non-Ricker, per-\(\sigma\), recoverable/sink, with the current
  public-demographics run kept as a labelled data-rich comparison arm;
  \item a short note of anything in the brief that conflicts with actual code.
\end{enumerate}

\end{document}
